\section{Conclusion}\label{sec:conclusion109}
Neural Bellman Operators remain a construction-and-learning approach to original recursive economic decisions. Their continuation maps must be judged by feasible decisions, original-objective losses and complete paid service work, rather than by an internal residual score alone. The centered tail-switch theorem identifies a specific reason that directional fitting can fail under nonlinear risk. Original-law support and strict completed-incumbent gates supply a verifiable route from a proposed action to an actual cost improvement. The recursive gain account and precision bound make explicit which uniform hypotheses and expenditures are needed beyond a finite successful catalogue.

In the original investment economy, the active-region Bellman compiler provides a strong semi-analytic control and a useful common backend. The executed catalogue confirms strict policy improvements but also their additional cost. It confirms exact equivalence between acquired ReLU geometry and conventional hinges, not a uniquely neural advantage. Conventional victories, adverse residual-fitting results and unresolved adoption regions remain part of the scientific record.

The next scientific tests are substantive: prospectively train and evaluate the tail-aware objective under independently verified margin information; establish a neural representation advantage against equally informed conventional geometry and resources; solve genuinely longer-horizon, multiple-control and multiple-innovation instances; and calibrate where certified policy gains justify acquisition and service expense. The Brownian recursive-utility bridge states one rigorous connection to NDU/FBSDE research without transferring evidence across different objectives. Neither the finite comparison nor successful typesetting completes these broader claims. The original economic problem, mathematical ambition and numerical-methods objective are retained.
